% Part: incompleteness
% Chapter: representability-in-q
% Section: minimization-representable

\documentclass[../../../include/open-logic-section]{subfiles}

\begin{document}

\olfileid{inc}{req}{min}
\olsection{التصغير المنتظم قابل للتمثيل في $\Th{Q}$}

لنتأمل البحث غير المحدود. افترض أن $g(x,z)$ منتظمة وقابلة للتمثيل في
$\Th{Q}$، ولتكن !!{formula}~$!A_g(x,z,y)$ ممثلة لها. ولتُعرّف $f$ بـ
$f(z)=\umin{x}{[g(x,z)=0]}$. نريد أن نجد !!a{formula}~$!A_f(z,y)$
تمثل~$f$. فقيمة~$f(z)$ هي ذلك العدد~$x$ الذي (أ) يحقق $g(x,z)=0$،
و(ب) هو أصغر عدد يحقق ذلك؛ أي إن $g(w,z)\neq0$ لكل $w<x$. ومن ثم يكون
الاختيار الآتي طبيعيًا:
\[
!A_f(z,y) \ident !A_g(y, z, \Obj 0) \land \lforall[w][(w < y \lif \lnot
  !A_g(w, z, \Obj 0))].
\]
وفي الحالة العامة ينبغي، بالطبع، أن نستبدل $z$ بـ$z_0$، \dots، $z_k$.

وسيتضمن البرهان مرة أخرى بعض اللمم عما تكون $\Th{Q}$ قوية بما يكفي
لإثباته.

\begin{lem}
\ollabel{lem:succ} لكل !!{constant}~$a$ ولكل عدد طبيعي~$n$،
\[
\Th{Q} \Proves \eq[(a' + \num n)][(a + \num n)'].
\]
\end{lem}

\begin{proof}
البرهان، كالمعتاد، بالاستقراء على~$n$. في الحالة الأساسية $n=0$،
يلزم أن نبيّن أن $\Th{Q}$ تثبت
$\eq[(a'+\Obj 0)][(a+\Obj 0)']$. ولدينا:
\begin{align}
  \Th{Q} & \Proves \eq[(a' + \Obj 0)][a'] \quad \text{بالبديهية $Q_4$}
  \ollabel{step1}\\
  \Th{Q} & \Proves  \eq[(a + \Obj 0)][a] \quad \text{بالبديهية $Q_4$}
  \ollabel{step2} \\
  \Th{Q} & \Proves \eq[(a + \Obj 0)'][a'] \quad \text{بـ\olref{step2}}
  \ollabel{step3} \\
  \Th{Q} & \Proves \eq[(a' + \Obj 0)][(a + \Obj 0)'] \quad
  \text{بـ\olref{step1} و\olref{step3}}\notag
\end{align}
وفي خطوة الاستقراء يجوز أن نفترض أننا أثبتنا
$\Th{Q}\Proves\eq[(a'+\num n)][(a+\num n)']$. ولما كانت
$\num{n+1}$ هي $\num n'$، فعلينا أن نبيّن أن $\Th{Q}$ تثبت
$\eq[(a'+\num n')][(a+\num n')']$. ولدينا:
\begin{align}
  \Th{Q} & \Proves \eq[(a' + \num n')][(a' + \num n)'] \quad
  \text{بالبديهية $!Q_5$} \ollabel{step5}\\
  \Th{Q} & \Proves \eq[(a' + \num n')][(a + \num n')'] \quad
  \text{بفرض الاستقراء والبديهية $!Q_5$ ومنطق المساواة} \ollabel{step6}\\
  \Th{Q} & \Proves \eq[(a' + \num n)'][(a + \num n')'] \quad
  \text{بـ\olref{step5} و\olref{step6}.} \notag
\end{align}
\end{proof}

ويجدر التنبيه مرة أخرى إلى أن هذا أضعف من القول إن $\Th{Q}$ تثبت
$\lforall[x][\lforall[y][(x'+y)=(x+y)']]$. فمع أن هذه
!!{sentence} صادقة في~$\Struct{N}$، لا تثبتها $\Th{Q}$.

\begin{lem}
\ollabel{lem:less-zero}
$\Th{Q} \Proves \lforall[x][\lnot x < \Obj 0]$.
\end{lem}

\begin{proof}
نقدّم البرهان بصورة غير صورية (أي لا نقدّم إلا إرشادات لبناء
!!{derivation} الصورية).

علينا أن نثبت $\lnot a<\Obj 0$ من أجل~$a$ اعتباطية. وبحسب تعريف
$<$، يلزم أن نثبت
$\lnot\lexists[y][\eq[(y'+a)][\Obj 0]]$ في $\Th{Q}$. سنفترض
$\lexists[y][\eq[(y'+a)][\Obj 0]]$ ونستخرج تناقضًا. افترض
$\eq[(b'+a)][\Obj 0]$. باستعمال $!Q_3$ نحصل على
$\eq[a][\Obj 0]\lor\lexists[y][\eq[a][y']]$. ونفصل حالتين.

الحالة 1: تصدق $\eq[a][\Obj 0]$. من $\eq[(b'+a)][\Obj 0]$ نحصل
على $\eq[(b'+\Obj 0)][\Obj 0]$. وبحسب البديهية~$!Q_4$ في
$\Th{Q}$، لدينا $\eq[(b'+\Obj 0)][b']$، ومن ثم
$\eq[b'][\Obj 0]$. لكن البديهية~$!Q_2$ تعطينا أيضًا
$\eq/[b'][\Obj 0]$، وهو تناقض.

الحالة 2: يوجد $c$ بحيث $\eq[a][c']$. عندئذ نحصل على
$\eq[(b'+c')][\Obj 0]$. وبحسب البديهية~$!Q_5$، نحصل على
$\eq[(b'+c)'][\Obj 0]$، فيناقض ذلك البديهية~$Q_2$ مرة أخرى.
\end{proof}

\begin{lem}
\ollabel{lem:less-nsucc}
لكل عدد طبيعي~$n$،
  \[
  \Th{Q} \Proves
  \lforall[x][(x < \num {n+1} \lif (\eq[x][\Obj 0] \lor \dots \lor
    \eq[x][\num n]))].
  \]
\end{lem}

\begin{proof}
نستعمل الاستقراء على~$n$. لننظر في الحالة الأساسية $n=0$. يلزم فيها
أن نبيّن $a<\num 1\lif\eq[a][\Obj 0]$ من أجل~$a$ اعتباطية. افترض
$a<\num 1$. عندئذ تعطينا البديهية المعرِّفة لـ$<$ الصيغة
$\lexists[y][\eq[(y'+a)][\Obj 0']]$ (لأن
$\num 1\ident\Obj 0'$).

افترض أن $b$ تحقق هذه الخاصية؛ أي لدينا
$\eq[(b'+a)][\Obj 0']$. علينا أن نبيّن $\eq[a][\Obj 0]$. بحسب
البديهية~$!Q_3$، إما $\eq[a][\Obj 0]$، وإما يوجد $c$ بحيث
$\eq[a][c']$. لا شيء نثبته في الحالة الأولى؛ فافترض
$\eq[a][c']$. عندئذ لدينا $\eq[(b'+c')][\Obj 0']$. وبحسب
البديهية~$!Q_5$ في $\Th{Q}$، لدينا
$\eq[(b'+c)'][\Obj 0']$. وبحسب البديهية~$!Q_1$، لدينا
$\eq[(b'+c)][\Obj 0]$. لكن هذا يعني، بحسب البديهية~$!Q_8$، أن
$c<\Obj 0$، وهو ما يناقض \olref{lem:less-zero}.

والآن إلى خطوة الاستقراء. نثبت حالة $n+1$ مفترضين حالة~$n$. افترض
$a<\num{n+2}$. وباستعمال $!Q_3$ مرة أخرى نفصل حالتين:
$\eq[a][\Obj 0]$، ووجود $b$ بحيث $\eq[a][b']$. في الحالة الأولى
ينتج
$\eq[a][\Obj 0]\lor\dots\lor\eq[a][\num{n+1}]$ بداهة. وفي
الحالة الثانية لدينا $b'<\num{n+2}$، أي
$b'<\num{n+1}'$. وبحسب البديهية~$!Q_8$ يوجد $c$ بحيث
$\eq[(c'+b')][\num{n+1}']$. وبحسب البديهية~$!Q_5$،
$\eq[(c'+b)'][\num{n+1}']$. وبحسب البديهية~$!Q_1$،
$\eq[(c'+b)][\num{n+1}]$، ولذلك $b<\num{n+1}$ بحسب
البديهية~$!Q_8$. وبفرض الاستقراء،
$\eq[b][\Obj 0]\lor\dots\lor\eq[b][\num n]$. ومن ذلك نحصل
بالمنطق على
$\eq[b'][\Obj 0']\lor\dots\lor\eq[b'][\num n']$، ومن ثم
$\eq[a][\num 1]\lor\dots\lor\eq[a][\num{n+1}]$ لأن
$\eq[a][b']$.
\end{proof}

\begin{lem}
  \ollabel{lem:trichotomy} لكل عدد طبيعي~$m$،
  \[
  \Th{Q} \Proves
  \lforall[y][((y < \num{m} \lor \num{m} < y) \lor \eq[y][\num{m}])].
  \]
\end{lem}

\begin{proof}
بالاستقراء على~$m$. تأمل أولًا الحالة $m=0$. بحسب~$!Q_3$،
$\Th{Q}\Proves
\lforall[y][(\eq[y][\Obj 0]\lor\lexists[z][\eq[y][z']])]$.
لتكن $a$ اعتباطية. فإما $\eq[a][\Obj 0]$، وإما يوجد~$b$ بحيث
$\eq[a][b']$. في الحالة الأولى لدينا أيضًا
$(a<\Obj 0\lor\Obj 0<a)\lor\eq[a][\Obj 0]$. أما إذا كانت
$\eq[a][b']$، فإن $\eq[(b'+\Obj 0)][(a+\Obj 0)]$ بحسب منطق
$\eq$. وبحسب $!Q_4$، $\eq[(a+\Obj 0)][a]$، ولذلك نحصل على
$\eq[(b'+\Obj 0)][a]$، ومن ثم
$\lexists[z][\eq[(z'+\Obj 0)][a]]$. وبحسب تعريف $<$ في $!Q_8$،
$\Obj 0<a$. وإذا كانت $\Obj 0<a$، كان أيضًا
$(\Obj 0<a\lor a<\Obj 0)\lor\eq[a][\Obj 0]$.

والآن افترض أن لدينا
\begin{align*}
  \Th{Q} & \Proves \lforall[y][((y < \num{m} \lor \num{m} < y) \lor
    \eq[y][\num{m}])]
  \intertext{ونريد أن نبيّن}
  \Th{Q} & \Proves \lforall[y][((y < \num{m+1} \lor \num{m+1} < y) \lor
    \eq[y][\num{m+1}])]
\end{align*}
لتكن $a$ اعتباطية. بحسب $!Q_3$، إما $\eq[a][\Obj 0]$، وإما يوجد
$b$ بحيث $\eq[a][b']$. في الحالة الأولى لدينا
$\eq[\num m'+a][\num{m+1}]$ بحسب $!Q_4$، ومن ثم
$a<\num{m+1}$ بحسب $!Q_8$.

تأمل الآن الحالة الثانية، $\eq[a][b']$. بحسب فرض الاستقراء،
$(b<\num m\lor\num m<b)\lor\eq[b][\num m]$.

الحد الأول $b<\num m$ يكافئ (بحسب $!Q_8$)
$\lexists[z][\eq[(z'+b)][\num m]]$. افترض أن $c$ تحقق هذه الخاصية.
فإذا كانت $\eq[(c'+b)][\num m]$، كانت أيضًا
$\eq[(c'+b)'][\num m']$. وبحسب $!Q_5$،
$\eq[(c'+b)'][(c'+b')]$. ومن ثم
$\eq[(c'+b')][\num m']$. ونحصل على
$\lexists[u][\eq[(u'+b')][\num{m+1}]]$ بالتعميم الوجودي على~$c$،
مع مراعاة أن $\num m'\ident\num{m+1}$. ولذلك إذا كان
$b<\num m$، كان $b'<\num{m+1}$، ومن ثم $a<\num{m+1}$.

وافترض الآن $\num m<b$، أي
$\lexists[z][\eq[(z'+\num m)][b]]$. ولتكن $c$ مثل هذا~$z$؛ أي
$\eq[(c'+\num m)][b]$. بالمنطق،
$\eq[(c'+\num m)'][b']$. وبحسب $!Q_5$،
$\eq[(c'+\num m')][b']$. ولما كانت $\eq[a][b']$ وكان
$\num m'\ident\num{m+1}$، نحصل على
$\eq[(c'+\num{m+1})][a]$. وبحسب $!Q_8$، $\num{m+1}<a$.

وأخيرًا، افترض $\eq[b][\num m]$. عندئذ، بالمنطق،
$\eq[b'][\num m']$، ومن ثم $\eq[a][\num{m+1}]$.

وهكذا نستطيع أن نستخرج من كل حد من فصل حالة~$m$ و$b$ الحد المناظر
لحالة~$m+1$ و$a$.
\end{proof}
  
\begin{prop}
\ollabel{prop:rep-minimization}
إذا كانت $!A_g(x,z,y)$ تمثل $g(x,z)$ في~$\Th{Q}$، فإن
\[
!A_f(z,y) \ident !A_g(y, z, \Obj 0) \land \lforall[w][(w < y \lif \lnot
  !A_g(w, z, \Obj 0))]
\]
تمثل $f(z)=\umin{x}{[g(x,z)=0]}$.
\end{prop}

\begin{proof}
نبيّن أولًا أنه إذا كان $f(n)=m$، فإن
$\Th{Q}\Proves !A_f(\num n,\num m)$؛ أي
\begin{align}
  \Th{Q} & \Proves !A_g(\num{m}, \num{n}, \Obj 0) \land \lforall[w][(w <
    \num{m} \lif \lnot !A_g(w, \num{n}, \Obj 0))]. \notag
  \intertext{ولما كانت $!A_g(x,z,y)$ تمثل $g(x,z)$ وكان $g(m,n)=0$
    إذا كان $f(n)=m$، فلدينا}
\Th{Q} & \Proves !A_g(\num{m}, \num{n}, \Obj 0). \notag
\intertext{وإذا كان $f(n)=m$، فإن $g(k,n)\neq0$ لكل $k<m$. ومن ثم}
\Th{Q} & \Proves \lnot !A_g(\num{k}, \num{n}, \Obj 0). \notag
\intertext{فنحصل على}
\Th{Q} & \Proves \lforall[w][(w < \num{m} \lif \lnot
  !A_g(w, \num{n}, \Obj 0))]. \ollabel{rep-less}
\end{align}
وذلك بـ\olref{lem:less-zero} إذا كان $m=0$، وبـ
\olref{lem:less-nsucc} فيما عدا ذلك.

والآن نبيّن أنه إذا كان $f(n)=m$، فإن
$\Th{Q}\Proves
\lforall[y][(!A_f(\num n,y)\lif\eq[y][\num m])]$. ونوجز الحجة
مرة أخرى بصورة غير صورية، تاركين صوغها الرسمي للقارئ.

افترض $!A_f(\num n,b)$. نستخرج منها (أ)
$!A_g(b,\num n,\Obj 0)$، و(ب)
$\lforall[w][(w<b\lif\lnot !A_g(w,\num n,\Obj 0))]$. وبحسب
\olref{lem:trichotomy}،
$(b<\num m\lor\num m<b)\lor\eq[b][\num m]$. وسنبيّن أن كلا من
$b<\num m$ و$\num m<b$ يؤدي إلى تناقض.

إذا كان $\num m<b$، فإن $\lnot !A_g(\num m,\num n,\Obj 0)$
ينتج من~(ب). لكن $m=f(n)$، ولذلك $g(m,n)=0$، ومن ثم
$\Th{Q}\Proves !A_g(\num m,\num n,\Obj 0)$ لأن $!A_g$ تمثل~$g$.
فنحصل على تناقض.

وافترض الآن $b<\num m$. لما كان
$\Th{Q}\Proves\lforall[w][(w<\num m\lif
\lnot !A_g(w,\num n,\Obj 0))]$ بحسب \olref{rep-less}، نحصل على
$\lnot !A_g(b,\num n,\Obj 0)$. وهذا يناقض~(أ) مرة أخرى.
\end{proof}

\end{document}
